\section{Preliminaries}
\label{sec:preliminaries}

A finite set $Z\subseteq\mathcal X$ is \emph{shattered} by $\mathcal H$ if
every map $b:Z\to\{-1,+1\}$ is the restriction of some member of
$\mathcal H$.  The VC dimension $\operatorname{VC}(\mathcal H)$ is the
largest cardinality of a shattered subset of $\mathcal X$.  This
setting-derived quantity is introduced only to eliminate the class cardinality
in the proof of the main theorem; it does not remain in the final bound.

All feature representations use the Euclidean inner product and the same
fixed map $\operatorname{sgn}_{\tau}$.  In particular, equality of two scores
preserves their predicted labels even when both scores are zero.  This
convention is used in the direct-sum construction proving the main theorem.
